% Pista ana-26j2-q4a · Anàlisi
% Procedència: PAU Catalunya 2026 · Convocatòria de juny · Sèrie 5 · Exercici 4 — Opció A
\documentclass[11pt,a4paper]{article}
\usepackage{pista}
\pistainfo{Catalunya 2026}{Convocatòria de juny}{Sèrie 5}

\begin{document}

\section*{\textcolor{blauPista}{Exercici 4 --- Opció A}}

\noindent D'una funció $f(x)$, sabem que passa pel punt $(-1, -1)$ i que coneixem la gràfica de la seva funció derivada $f'(x)$ (amb punts de tall amb l'eix $OX$ a $x = -2$, $x = 0$ i $x = 1$).

\subsection*{4.1. \normalfont Calculeu l'equació de la recta tangent a la gràfica de la funció $f(x)$ en el punt $(-1, -1)$. \hfill\textnormal{[1 punt]}}

\pista{Per a l'equació de la recta tangent en un punt, necessites el punt (que ja et donen, $(-1, -1)$) i el pendent, que és $f'(-1)$. Aquest valor no l'has de calcular: llegeix-lo directament de la gràfica de $f'(x)$ que et donen (mira quin valor pren la corba a $x = -1$).}

\subsection*{4.2. \normalfont Estudieu si $f(x)$ té punts crítics (punts on $f'(x) = 0$) i, en cas que en tingui, determineu-ne les abscisses i classifiqueu-los segons si són màxims, mínims o punts d'inflexió de $f(x)$. \hfill\textnormal{[1 punt]}}

\pista{Els punts crítics de $f(x)$ són, per definició, els punts on $f'(x) = 0$: identifica'ls directament sobre la gràfica (els talls de la corba $f'(x)$ amb l'eix $OX$). Per classificar-los com a màxims, mínims o punts d'inflexió, no cal calcular res més: n'hi ha prou d'observar el signe de $f'(x)$ (és a dir, si la gràfica de $f'$ està per damunt o per sota de l'eix $OX$) abans i després de cada punt crític. Compte especial amb un d'ells: si el signe de $f'$ no canvia (és positiva --o negativa-- abans i després), aleshores no és ni un màxim ni un mínim, sinó un punt d'inflexió de $f(x)$.}

\subsection*{4.3. \normalfont Sabent, a més, que $f(0) = -0{,}2$, calculeu l'àrea entre la gràfica de $f'(x)$ i l'eix de les abscisses, des de $x = -1$ fins a $x = 0$. \hfill\textnormal{[0,5 punts]}}

\pista{Recorda la regla de Barrow: com que $f(x)$ és, per definició, una primitiva de $f'(x)$, la integral definida de $f'(x)$ entre $-1$ i $0$ es calcula com $f(0) - f(-1)$, sense necessitat de trobar l'expressió explícita de $f'(x)$ ni integrar res. Ja tens els dos valors que necessites: un te'l donen a l'enunciat, i l'altre el coneixies des del principi (perquè és el mateix punt de l'apartat 4.1). Comprova, mirant la gràfica de $f'(x)$, si aquesta canvia de signe en aquest interval abans de donar el resultat final com a àrea.}

\end{document}
